\documentclass[tikz,border=4pt]{standalone}
% Paleta y fuentes compartidas por el libro y las figuras (única fuente de verdad).
\RequirePackage{xcolor}
\definecolor{oxido}{HTML}{B7410E}
\definecolor{cobre}{HTML}{D27D2D}
\definecolor{carbon}{HTML}{2B2B2B}
\definecolor{azul}{HTML}{1F5F8B}
\definecolor{arena}{HTML}{F4EDE4}

\usepackage{fontspec}
\setsansfont{FiraSans}[Extension=.otf,UprightFont=*-Regular,BoldFont=*-SemiBold,ItalicFont=*-Italic,BoldItalicFont=*-SemiBoldItalic]
\setmonofont{FiraMono}[Extension=.otf,UprightFont=*-Regular,BoldFont=*-Medium,Scale=0.95]
\usetikzlibrary{arrows.meta,calc,positioning}
\begin{document}
\sffamily\small
% Línea de tiempo vertical: año a la izquierda, hito a la derecha.
\begin{tikzpicture}[
  anio/.style={font=\sffamily\bfseries\small,text=oxido,anchor=east},
  hito/.style={font=\sffamily\small,text=carbon,anchor=west,text width=96mm,align=left},
  punto/.style={circle,fill=oxido,inner sep=2.1pt},
  claro/.style={circle,draw=oxido,fill=white,line width=1pt,inner sep=1.8pt}]
  \def\paso{0.82}
  \draw[carbon!25,line width=1.2pt] (0,0.3) -- (0,-13*\paso-0.3);
  \foreach \a/\t/\e [count=\i from 0] in {
    2006/{Graydon Hoare empieza Rust como proyecto personal}/claro,
    2009/{Mozilla patrocina el proyecto}/claro,
    2010/{Presentación en el Mozilla Summit}/claro,
    2011/{\texttt{rustc} se compila a sí mismo}/claro,
    2012/{Rust 0.1; nace Servo en Mozilla Research}/claro,
    2015/{\textbf{Rust 1.0}: compromiso de estabilidad}/punto,
    2016/{Primer año como lenguaje «más amado» en Stack Overflow}/claro,
    2017/{Stylo, escrito en Rust, llega a Firefox 57}/claro,
    2018/{Edición 2018 (Rust 1.31)}/claro,
    2020/{Despidos en Mozilla; Servo pasa a la Linux Foundation}/claro,
    2021/{\textbf{Nace la Rust Foundation}}/punto,
    2022/{Linux 6.1 incluye soporte inicial para Rust}/claro,
    2023/{Leadership Council; primer código Rust en el núcleo de Windows}/claro,
    2025/{Edición 2024 (Rust 1.85); Rust deja de ser experimental en Linux}/punto}
  {
    \node[\e] at (0,-\i*\paso) {};
    \node[anio] at (-0.35,-\i*\paso) {\a};
    \node[hito] at (0.3,-\i*\paso) {\hyphenpenalty=10000 \t};
  }
\end{tikzpicture}
\end{document}
